
前面的 $q\bar q\to gg$ 已经说明：一旦外态含胶子，非阿贝尔三胶子顶角不再是可有可无的修正，而是 Ward 恒等式所要求的树级组成部分。要把强子对撞机中最常见的 Born 硬核补齐，还需要再看两个含初态胶子的通道 $qg\to qg$ 与 $gg\to gg$。它们的意义不只是“再多算两个过程”：前者同时包含基本表示与伴随表示的颜色代数，后者则完全由胶子自相互作用组成，因而把 $C_F$ 与 $C_A$ 在可观测散射核中的不同角色直接显露出来。

先保留夸克质量 $m_q$，取
\begin{equation*}
 q(p_1,i)+g(k_1,a)\to q(p_2,j)+g(k_2,b),
 \qquad
 p_1^2=p_2^2=m_q^2c^2,
 \qquad k_1^2=k_2^2=0,
\end{equation*}
并定义
\begin{equation*}
 s=(p_1+k_1)^2,
 \qquad
 t=(p_1-p_2)^2,
 \qquad
 u=(p_1-k_2)^2,
 \qquad s+t+u=2m_q^2c^2.
\end{equation*}
树级振幅由 $s$、$u$ 道夸克交换以及 $t$ 道胶子交换三部分组成。把外胶子偏振矢量记为 $\varepsilon_\mu(k_1)$ 与 $\varepsilon_\nu^*(k_2)$，则可写成
\begin{equation}
{
\begin{aligned}
\mathcal M_{qg}(m_q)
=g_s^2\,\bar u(p_2)\Bigg[&
 T^bT^a\,\gamma^\nu
 \frac{\slashed p_1+\slashed k_1+m_qc}{s-m_q^2c^2}\gamma^\mu
 +T^aT^b\,\gamma^\mu
 \frac{\slashed p_1-\slashed k_2+m_qc}{u-m_q^2c^2}\gamma^\nu\\
&+\frac{\ii f^{abc}T^c}{t}\,
 \Gamma^{\rho\mu\nu}(p_1-p_2,k_1,-k_2)\gamma_\rho
\Bigg]u(p_1)\,
\varepsilon_\mu(k_1)\varepsilon_\nu^*(k_2).
\end{aligned}}
\label{eq:qcd-qg-amplitude-dp94}
\end{equation}
其中 $\Gamma^{\rho\mu\nu}$ 是前文已经由 $F^a_{\mu\nu}F^{a\mu\nu}$ 推出的三胶子顶角张量。第一条腿属于 $t$ 道内部胶子；因为三胶子规则把三条动量都定义为流入顶角，所以它携带的是 $p_1-p_2=k_2-k_1$，并满足 $(p_1-p_2)+k_1-k_2=0$。这条动量和是使用顶角公式前必须先做的检查。三项的相对系数由同一个协变导数与场强张量固定；若去掉第三项，纵向偏振替换不会相消。

有限 $m_q$ 的振幅由\eqnrefs{eq:qcd-qg-amplitude-dp94} 完整定义。强子对撞机的轻 parton Born 核进一步取严格连续极限 $m_q\to0$；此时 $s+t+u=0$。对初态夸克的 $N_c$ 个颜色和两个自旋、初态胶子的 $N_c^2-1$ 个颜色和两个物理偏振平均，并对末态求和，得到无质量部分子极限的闭式
\begin{equation}
{
\overline{|\mathcal M(qg\to qg)|^2}
=g_s^4(s^2+u^2)
\left[
 \frac1{t^2}-\frac{C_F}{C_A}\frac1{su}
\right],
\qquad
C_F=\frac{N_c^2-1}{2N_c},\quad C_A=N_c.}
\label{eq:qcd-qg-M2-dp94}
\end{equation}
这里第一项是前向 $t$ 道胶子交换产生的强增强，第二项来自夸克交换与非阿贝尔颜色排序之间的相干组合。对 QCD 的 $N_c=3$，$C_F/C_A=4/9$，恢复常用的
\begin{equation*}
 \overline{|\mathcal M|^2}
 =g_s^4(s^2+u^2)
 \left(\frac1{t^2}-\frac{4}{9su}\right).
\end{equation*}
由于物理区 $s>0$、$t<0$、$u<0$，第二项实际上为正；写成 $-1/(su)$ 的形式是为了保持交叉对称的 Mandelstam 结构。

纯胶子散射
\begin{equation*}
 g(k_1,a)+g(k_2,b)\to g(k_3,c)+g(k_4,d)
\end{equation*}
包含 $s,t,u$ 三个三胶子交换图以及一个四胶子接触图。用 Jacobi 恒等式把接触项分配到三个交换道后，完整树级振幅可按伴随表示颜色张量写成
\begin{equation}
{
 \mathcal M^{abcd}
 =g_s^2\left[
 f^{abe}f^{cde}\mathcal K_s
 +f^{ace}f^{bde}\mathcal K_t
 +f^{ade}f^{bce}\mathcal K_u
 \right],}
\label{eq:qcd-gg-color-amplitude-dp94}
\end{equation}
其中 $\mathcal K_s,\mathcal K_t,\mathcal K_u$ 包含相应交换图与四胶子接触项的洛伦兹结构。单独任何一项都依赖规范选择；只有它们的完整和在外腿上满足 Ward 恒等式。对一般 $SU(N_c)$ 完成颜色、偏振求和与初态平均后得到
\begin{equation}
{
\overline{|\mathcal M(gg\to gg)|^2}
=\frac{4g_s^4C_A^2}{N_c^2-1}
\left(
3-\frac{su}{t^2}-\frac{st}{u^2}-\frac{tu}{s^2}
\right).}
\label{eq:qcd-gg-M2-dp94}
\end{equation}
对 $N_c=3$，整体颜色系数成为 $9/2$。与 $qq'$ 或 $qg$ 相比，\eqnrefs{eq:qcd-gg-M2-dp94} 同时在 $t\to0$ 与 $u\to0$ 出现交换增强，反映了两个完全相同的无质量规范玻色子既可沿前向也可沿后向交换。式中的 $C_A^2$ 还直接表明：胶子自相互作用的强度由伴随表示 Casimir 控制，而不是由夸克的 $C_F$ 控制。

对于任意无质量 $2\to2$ Born 过程，统一相空间关系为
\begin{equation}
 \frac{\dd\sigma}{\dd t}
 =\frac{\hbar^2}{16\pi s^2}\overline{|\mathcal M|^2}.
 \label{eq:qcd-partonic-dsdt-master-dp94}
\end{equation}
因此 \eqnrefs{eq:qcd-qg-M2-dp94,eq:qcd-gg-M2-dp94} 已经足以直接生成角分布。\figref{fig:qcd-qg-gg-angular-dp94} 把 $N_c=3$ 时的无量纲 Born 核除以各自在 $90^\circ$ 的值；这样既去掉共同的 $g_s^4$ 与 $1/s$ 尺度，又保留由交换道和颜色结构决定的角形状。

\begin{figure}[htbp]
 \centering
 \includegraphics[width=.74\textwidth]{qcd_qg_gg_angular_dp94.pdf}
 \caption{QCD 中 $qg\to qg$ 与 $gg\to gg$ 的树级无量纲角分布核，均按 $90^\circ$ 归一化。$qg$ 通道主要显示 $t$ 道前向增强；$gg$ 通道还因相同外胶子而出现明显的后向交换增强。曲线由\eqnrefs{eq:qcd-qg-M2-dp94,eq:qcd-gg-M2-dp94} 逐点计算。}
 \label{fig:qcd-qg-gg-angular-dp94}
\end{figure}

